\chapter{TỔNG KẾT}
\label{ch:tong_ket}

\section{Tóm lược kết quả đạt được}

Đề tài đã khảo sát hai kỹ thuật nén mô hình --- chưng cất tri thức và lượng tử hoá sau huấn luyện
--- cho bài toán phát hiện biển báo giao thông, trên bộ dữ liệu \textit{pkdarabi/cardetection} gồm
4.969 ảnh với 6.012 hộp bao thuộc 15 lớp. Năm cấu hình được huấn luyện và đánh giá trong cùng một
lượt chạy, bằng cùng một bộ tính độ đo.

Đối chiếu với ba câu hỏi nghiên cứu đặt ra ở Chương~\ref{ch:gioi_thieu}:

\textbf{Câu hỏi 1.1} về chưng cất tri thức nhận câu trả lời \emph{phủ định}. Mô hình học trò được
chưng cất đạt mAP@0,5 bằng 0,9014, thấp hơn chính nó khi huấn luyện thông thường (0,9348) 3,34 điểm
phần trăm. Cơ chế chưng cất đã được xác minh là thực sự hoạt động thông qua cột mất mát riêng trong
nhật ký huấn luyện, giảm từ 12,1663 xuống 1,1608. Nguyên nhân khả dĩ nhất là khoảng cách năng lực
giữa thầy và trò quá hẹp --- chỉ 1,72 điểm phần trăm --- khiến lượng tri thức có thể truyền đạt
không bù nổi chi phí nhiễu.

\textbf{Câu hỏi 1.2} về lượng tử hoá cho thấy một sự đánh đổi phức tạp hơn nhiều so với thường được
trình bày. Về dung lượng, INT8 đạt mục tiêu rõ ràng: giảm 70,3\%. Về tốc độ, kết quả trái ngược kỳ
vọng --- bản INT8 chậm hơn bản FP32 hơn hai lần do môi trường thực thi không có nhân tính toán số
nguyên được tối ưu. Về độ chính xác, mAP@0,5 chỉ giảm 0,0038 nhưng $AP$ trên vật thể nhỏ giảm
0,0247.

\textbf{Câu hỏi 1.3} về tính đại diện của mAP tổng là câu hỏi thu được câu trả lời có giá trị nhất.
Mức suy giảm trên vật thể nhỏ lớn gấp 1,85 lần mAP tổng ở phép chưng cất, và gấp 6,5 lần ở phép
lượng tử hoá. Nguyên nhân nằm ở phân bố dữ liệu: 51,30\% số hộp bao thuộc nhóm lớn --- nhóm mà mọi
cấu hình đều xử lý tốt như nhau --- nên chúng chi phối giá trị trung bình và pha loãng suy giảm ở
nhóm khó.

Khi loại 65 ảnh bị rò rỉ và đo lại trên 573 ảnh sạch (Mục~\ref{sec:tap_sach}), tỉ số ở phép chưng
cất tăng từ 1,85 lên \textbf{2,66} lần. Kết luận vì vậy không phải tạo tác của dữ liệu bẩn --- làm
sạch dữ liệu chỉ khiến nó rõ hơn.

\section{Đóng góp}

Ba đóng góp chính có thể rút ra.

Thứ nhất, về mặt phương pháp thực nghiệm, đề tài thiết lập một thiết kế có đối chứng đúng nghĩa: mô
hình học trò được huấn luyện hai lần với cùng số chu kỳ, cùng hạt ngẫu nhiên, cùng kích thước ảnh,
chỉ khác ở việc bật hay tắt chưng cất. Kèm theo đó là năm ngưỡng chấp nhận khai báo trước khi chạy,
và cơ chế xác minh rằng kỹ thuật khảo sát đã thực sự được kích hoạt.

Thứ hai, và quan trọng nhất, đề tài cung cấp bằng chứng định lượng cho thấy \textbf{mAP tổng che
giấu chi phí thật của việc nén mô hình}. Một cấu hình được tổng kết là ``giảm 70\% dung lượng, chỉ
mất 0,4\% mAP'' thực chất mất 2,47 điểm phần trăm ở nhóm vật thể nhỏ. Với bài toán biển báo giao
thông, nơi biển ở xa quyết định thời gian phản ứng của xe, sự khác biệt này có ý nghĩa an toàn trực
tiếp. Kiến nghị rút ra là mọi báo cáo về nén mô hình cho bài toán phát hiện nên kèm $AP$ tách theo
nhóm kích thước --- chi phí bổ sung gần như bằng không nếu đã dùng bộ đánh giá COCO.

Ngoài ra, luận điểm này được kiểm chứng độc lập trên dữ liệu ngoài phân bố bằng \emph{chính} mô
hình học trò được chưng cất: trên ba đoạn video đường phố, tỉ lệ tăng số phát hiện giảm đơn điệu
theo kích thước vật thể --- 13,0 lần ở nhóm rất nhỏ, 4,3 lần ở nhóm nhỏ, 2,3 lần ở nhóm vừa, và
\emph{giảm} một nửa ở nhóm lớn. Gradient sạch như vậy là chữ ký của cơ chế đã nêu, chứ không phải
của nhiễu ngẫu nhiên.

Thứ ba, đề tài chỉ ra hai vấn đề của một bộ dữ liệu công khai đang được dùng rộng rãi: 38,68\% số
ảnh là ảnh cắt cận cảnh kiểu GTSRB khiến phân bố kích thước bị lưỡng cực, và 10,2\% tập kiểm tra có
bản sao trong tập huấn luyện.

\section{Hạn chế}

Đề tài có những hạn chế sau, được nêu đầy đủ để người đọc đánh giá đúng phạm vi giá trị của kết
quả.

\textbf{Kích thước lô không đồng nhất giữa hai nhánh học trò.} Nhánh chưng cất dùng lô 8, nhánh đối
chứng dùng 24, do ràng buộc bộ nhớ card đồ hoạ khi phải nạp đồng thời cả hai mô hình. Đây là yếu tố
gây nhiễu chưa kiểm soát được, và vì nó, kết luận đúng đắn là ``trong cấu hình cụ thể này, chưng
cất không mang lại lợi ích'' chứ không phải ``chưng cất có hại''.

\textbf{Hệ số chưng cất chưa được quét.} Giá trị 6,0 là mặc định của thư viện. Không loại trừ khả
năng một giá trị nhỏ hơn sẽ cho kết quả khác.

\textbf{Bộ dữ liệu có rò rỉ giữa các tập.} 65 trong số 638 ảnh kiểm tra có bản sao trong tập huấn
luyện, tức 10,2\%. Bộ chỉ số đã được đo lại trên 573 ảnh sạch nên hạn chế này không còn ảnh hưởng
tới kết luận, nhưng vẫn cần nêu vì các bảng chính của báo cáo vẫn báo cáo trên tập đầy đủ để so sánh
được với các công bố khác dùng cùng bộ dữ liệu. Cần lưu ý thêm rằng rò rỉ \emph{không} phân bố đều:
71,3\% số hộp bao trong phần rò rỉ thuộc nhóm vật thể nhỏ.

\textbf{Chỉ khảo sát một cặp thầy--trò.} Kết luận về chưng cất được rút ra từ đúng một cặp
YOLO26s--YOLO26n. Muốn kết luận tổng quát cần khảo sát nhiều cặp có khoảng cách năng lực khác nhau.

\textbf{Lượng tử hoá chỉ đo trên một môi trường thực thi.} Kết quả tốc độ của INT8 phụ thuộc mạnh
vào nhân tính toán của phần cứng đích, nên con số thu được ở đây không nên tổng quát hoá cho các
môi trường khác.

\textbf{Tập lớp bị giới hạn.} Mười lăm lớp chỉ gồm đèn tín hiệu và biển giới hạn tốc độ. Các nhóm
biển cấm, biển nguy hiểm và biển chỉ dẫn hoàn toàn không có trong tập nhãn, nên mô hình không thể
phát hiện chúng dù chất lượng huấn luyện có tốt đến đâu.

\textbf{Phép kiểm chứng trên video không có nhãn thật.} Ba đoạn video ở Mục~\ref{sec:video} chưa
được gán nhãn, nên mục đó chỉ báo cáo các đại lượng quan sát được chứ không báo cáo độ chính xác.
Không thể loại trừ hoàn toàn khả năng một phần phát hiện thêm là dương tính giả, dù gradient đơn
điệu theo kích thước nghiêng mạnh về giả thuyết ngược lại.

\textbf{Phép so sánh INT8 với FP32 chưa được tái lập độc lập.} Ba cấu hình PyTorch đã được đo lại
tại chỗ và khớp đến bốn chữ số thập phân với sổ tay gốc. Hai bản ONNX thì chưa: bản FP32 không có
sẵn ở môi trường đo lại, còn bản INT8 cho kết quả lệch 0,0066 do khác nhân tính toán. Kết luận về
lượng tử hoá vì vậy vẫn dựa trên một lần đo duy nhất, khác với kết luận về chưng cất.

\section{Hướng phát triển}

Các hướng phát triển được sắp theo thứ tự ưu tiên, từ chi phí thấp tới chi phí cao.

\textbf{Ưu tiên 1 --- Tái lập phép đo ONNX.} Tải lại bản FP32 và đo cả hai bản ONNX trong cùng một
môi trường thực thi, để phép so sánh INT8 với FP32 đạt mức tin cậy ngang với phần chưng cất. Chi phí
gần như bằng không vì không phải huấn luyện lại.

\textbf{Ưu tiên 2 --- Chạy lại nhánh chưng cất với kích thước lô đồng nhất.} Dùng tích luỹ gradient
để đạt kích thước lô hiệu dụng bằng 24 như nhánh đối chứng, qua đó loại bỏ yếu tố gây nhiễu lớn
nhất hiện nay. Chỉ sau bước này mới có thể kết luận chắc chắn về tác dụng của chưng cất.

\textbf{Ưu tiên 3 --- Quét hệ số chưng cất.} Thử các giá trị nhỏ hơn 6,0 để kiểm chứng giả thuyết
rằng thành phần bắt chước đang lấn át tín hiệu từ nhãn cứng.

\textbf{Ưu tiên 4 --- Mở rộng khoảng cách thầy--trò.} Dùng một mô hình thầy lớn hơn hẳn để tạo
khoảng cách năng lực đủ rộng, qua đó kiểm chứng trực tiếp giả thuyết về điều kiện áp dụng của chưng
cất.

\textbf{Ưu tiên 5 --- Làm sạch và cân bằng lại bộ dữ liệu.} Loại bỏ các ảnh trùng, sau đó thử huấn
luyện lại khi đã loại hoặc giảm mạnh nhóm 1.922 ảnh cắt kiểu GTSRB. Giả thuyết là mAP tổng sẽ giảm
nhưng $AP$ trên nhóm vật thể nhỏ sẽ tăng --- bản thân việc kiểm chứng giả thuyết này đã là một thí
nghiệm có giá trị.

\textbf{Ưu tiên 6 --- Gán nhãn một tập khung hình video và đo $AP$ thật.} Phép kiểm chứng ở
Mục~\ref{sec:video} mới dừng ở mức nhất quán với giả thuyết. Gán nhãn thủ công vài trăm khung hình
từ ba video đã có sẽ biến nó thành một phép đo $AP$ đúng nghĩa trên dữ liệu ngoài phân bố --- thứ có
giá trị hơn hẳn mọi chỉ số trên tập kiểm tra cùng phân bố.

\textbf{Ưu tiên 7 --- Điều chỉnh kiến trúc cho vật thể nhỏ.} Vì toàn bộ phân tích ở
Chương~\ref{ch:ket_qua} đều chỉ về nhóm vật thể nhỏ, hướng cải thiện tự nhiên là bổ sung đầu dự
đoán P2 hoạt động ở bước sải 4 thay vì 8, hoặc huấn luyện ở kích thước ảnh lớn hơn. Đây là biện
pháp rẻ hơn nhiều so với chuyển sang một mô hình lớn hơn.

\textbf{Ưu tiên 8 --- Bổ sung dữ liệu đúng miền.} Không kỹ thuật nào thay thế được dữ liệu thật
đúng miền. Nếu cần bổ sung dữ liệu công khai, TT100K~\cite{zhu2016tt100k} và Mapillary Traffic
Sign~\cite{ertler2020mapillary} là hai lựa chọn phù hợp vì đều là ảnh cảnh đường thật, qua đó lấp
được đúng vùng kích thước trung gian mà bộ dữ liệu hiện tại đang thiếu.
